% 本文件作为当前原理章节的入口；原有 content/2principle.tex 保持不变。
\section{实验原理}

%  本节参考课程讲义《密码分析的自动化搜索模型》中关于非零位置点集的建模方法。\footnote{胡凯：《密码分析的自动化搜索模型》，2026年10月20日，Part 1，第10--27页；线性掩码传播参见第45--46页。}讲义以三位S盒$\chi_3$说明方法，本实验则按实验指导书使用四位CipherFour S盒；两者的建模步骤相同，但点的维数、具体系数和实验结果需要分别确定。

\subsection{DDT与LAT}
根据课程所学的差分分析和线性分析，这里简单复述一下S盒的差分分布表(DDT)和线性分布表(LAT)。

\subsubsection{DDT}
对于一个$m \times n$的S盒，即输入$m$比特、输出$n$比特的映射$S:\mathbb{F}_2^m \to \mathbb{F}_2^n$，记满足输入差分$\alpha$、输出差分$\beta$的输入集合为
$$
IN_S(\alpha ,\beta) = \{X \in \mathbb{F}_2^m :S(X) \oplus S(X \oplus \alpha) = \beta\}.
$$
DDT中的表项为该集合的元素个数，即
$$
\operatorname{DDT}[\alpha][\beta] = \# IN_S(\alpha,\beta).
$$
可见这是一张$2^m \times 2^n$的表，构建时可以先遍历输入差分$\alpha \in \mathbb{F}_2^m$共$2^m$种可能，内层循环遍历输入$X \in \mathbb{F}_2^m$共$2^m$种可能，根据计算出的输出差分更新表项。因此，构建DDT需要检查$2^m \times 2^m = 2^{2m}$组$(\alpha,X)$对应的差分传播。

\subsubsection{LAT}

对于同一个$m \times n$的S盒，LAT是一个$2^m \times 2^n$的表，其中第$\alpha$行第$\beta$列的表项为
$$
\operatorname{LAT}[\alpha][\beta]
=\#\{x\in\mathbb{F}_2^m\mid \alpha\cdot x\oplus\beta\cdot S(x)=0\}-2^{m-1}.
$$
这里$\cdot$表示二进制向量的模$2$内积。该表项是输入、输出掩码为$\alpha,\beta$时，使线性表达式$\alpha\cdot x\oplus\beta\cdot S(x)$等于$0$的输入个数与$2^{m-1}$的差，其值与输入总数$2^m$的比值恰好为线性近似的偏差：
$$
\epsilon(\alpha,\beta)
=\Pr[\alpha\cdot x=\beta\cdot S(x)]-\frac{1}{2}
=\frac{\operatorname{LAT}[\alpha][\beta]}{2^m}.
$$
其中概率是在输入$x$均匀分布于$\mathbb{F}_2^m$时计算的。构建该表时，可以先遍历输入掩码$\alpha\in\mathbb{F}_2^m$，内层遍历输出掩码$\beta\in\mathbb{F}_2^n$，最内层遍历输入$x\in\mathbb{F}_2^m$，统计线性表达式为$0$的输入个数。因此，共需要检查$2^m\times2^n\times2^m=2^{2m+n}$组$(\alpha,\beta,x)$对应的线性表达式。


\subsection{S盒传播的统一建模约定}
\label{subsec:lab4-modeling-conventions}

本实验使用$4$比特输入、$4$比特输出的S盒。对于输入、输出差分对或掩码对$(a,b)$，将两者分别按最高位到最低位展开，再按输入在前、输出在后的顺序拼接为
\[
\operatorname{enc}(a,b)
=(a_3,a_2,a_1,a_0,b_3,b_2,b_1,b_0)\in\{0,1\}^{8},
\]
其中$a=\sum_{i=0}^{3}2^i a_i$，$b=\sum_{i=0}^{3}2^i b_i$。据此，分别定义差分传播和线性掩码传播的合法点集为
\begin{align*}
V_{\mathrm D}
&=\{\operatorname{enc}(a,b)\mid
  a,b\in\{0,\ldots,15\},\ \operatorname{DDT}[a][b]>0\},\\
V_{\mathrm L}
&=\{\operatorname{enc}(a,b)\mid
  a,b\in\{0,\ldots,15\},\ \operatorname{LAT}[a][b]\ne0\}.
\end{align*}

本报告的DDT与LAT均按输入在行、输出在列存储；讲义中的传播矩阵按输出在行、输入在列记作$M[b,a]$或$C[b,a]$。对于本实验的四位S盒，两种记法的对应关系为
\[
M[b,a]=\frac{\operatorname{DDT}[a][b]}{16},
\qquad
C[b,a]=\frac{\operatorname{LAT}[a][b]}{8}.
\]
其中$C$为相关度矩阵，LAT沿用前文的计数减半定义。矩阵转置的存储约定和非零缩放均不改变同一对$(a,b)$是否属于合法集合，但提取点集时必须保持输入、输出顺序一致。

LAT中的负数表示负偏差，仍属于非零相关的合法掩码对，因此构造$V_{\mathrm L}$时不能只保留正数表项。两个集合均保留$(a,b)=(0,0)$对应的全零向量，以描述未活跃S盒的传播。

下面统一用$V$表示$V_{\mathrm D}$或$V_{\mathrm L}$，用$z=(z_1,\ldots,z_8)$表示待约束的二进制变量。建模要求是：约束在这八个原始变量上的可行点集恰好为$V$；如果引入辅助变量，则取全部可行解在$z$上的投影。四种方法分别对两个集合应用相同的构造规则，仅刻画传播的可行性，不额外编码具体的概率或相关度权重。

\subsection{one-hot法}
\label{subsec:lab4-one-hot}

\paragraph{基本原理}

One-hot（独热编码）法按讲义中的三个步骤建立模型：先列出合法点并为每个点引入一个选择变量，再要求恰好选中一个点，最后按坐标恢复被选中点的值。该方法直接使用点集，不需要先计算凸包或进行逻辑化简。

\paragraph{数学建模}

沿用第\ref{subsec:lab4-modeling-conventions}节的统一记号，将合法点集写为
\[
V=\{v^{(1)},\ldots,v^{(N)}\},
\qquad
v^{(i)}=(v^{(i)}_1,\ldots,v^{(i)}_8),
\]
其中$N=|V|$。将各点作为列向量排成常数矩阵$A=(v^{(1)},\ldots,v^{(N)})$，并令$\lambda=(\lambda_1,\ldots,\lambda_N)^{\mathsf T}$，则唯一选择和坐标恢复分别写为
\[
\boldsymbol{1}^{\mathsf T}\lambda=1,
\qquad
z=A\lambda.
\]
逐坐标展开，并明确原始变量与选择变量的二进制取值，得到模型
\begin{equation}
\boxed{
\begin{aligned}
&\sum_{i=1}^{N}\lambda_i=1,\\
&z_j=\sum_{i=1}^{N}v^{(i)}_j\lambda_i,
&&j=1,\ldots,8,\\
&\lambda_i\in\{0,1\},
&&i=1,\ldots,N,\\
&z_j\in\{0,1\},
&&j=1,\ldots,8.
\end{aligned}
}
\label{eq:lab4-one-hot-model}
\end{equation}

其中，第一项约束保证唯一选择，第二项约束负责逐坐标恢复该合法点。由于$v^{(i)}_j$均为已知的$0/1$常数，第$j$个坐标也可稀疏地写成
\[
z_j=\sum_{i:\,v^{(i)}_j=1}\lambda_i.
\]
因此，程序只需累加在该坐标取$1$的点对应的选择变量；若没有这样的点，则该坐标等于$0$。所有等式都不含变量之间的乘积，仍为线性约束。模型中的原始变量与辅助选择变量应分别统计，不能把选择变量数当作模型的总变量数。

\paragraph{正确性分析}

对于模型的任意可行解，由$\lambda_i$的二进制取值及其和为$1$可知，存在唯一的$k$使得$\lambda_k=1$，其余选择变量均为$0$。代入坐标等式，得到$z_j=v^{(k)}_j$，从而$z=v^{(k)}\in V$。因此，模型不会引入非法点。

反过来，任取合法点$v^{(k)}\in V$，令$z=v^{(k)}$、$\lambda_k=1$，其余$\lambda_i=0$，即可满足全部约束。因此，每个合法点都能由模型表示。综上，
\[
\left\{z\in\{0,1\}^{8}\mid
\exists\lambda\in\{0,1\}^{N}\text{满足式}\eqref{eq:lab4-one-hot-model}
\right\}=V.
\]

\paragraph{用于差分与线性传播}

差分建模时取$V=V_{\mathrm D}$，每个选择变量对应一个DDT非零位置；线性建模时取$V=V_{\mathrm L}$，每个选择变量对应一个LAT非零位置。也可以直接用$(a,b)$给选择变量编号，记作$\lambda_{(a,b)}$，从而与原表中的位置一一对应。两类模型分别建立各自的选择变量，坐标约束形式保持一致；全零点的选择变量也需要保留，以表示未活跃S盒。表项的具体数值不作为这里的优化目标，枚举结果应覆盖全部合法点。

\subsection{逐点排除法}
\label{subsec:lab4-point-exclusion}

\paragraph{基本原理}

逐点排除法从所有八位二进制向量出发，对每个非法点建立一条只排除该点的不等式。将这些不等式同时加入模型后，所有非法点都被排除，而合法点全部保留。

\paragraph{数学建模}

设非法点集为
\[
U=\{0,1\}^{8}\setminus V.
\]
对于任意$u=(u_1,\ldots,u_8)\in U$，定义$z$与$u$的Hamming距离
\[
d_u(z)
=\sum_{j:u_j=0}z_j+\sum_{j:u_j=1}(1-z_j).
\]
当$u_j=0$时，$z_j$恰好表示该位是否不同；当$u_j=1$时，$1-z_j$恰好表示该位是否不同。因此，在二进制变量上，$d_u(z)$等于两向量不同坐标的数量。

要求$z$至少有一位不同于每个非法点，可得模型
\[
\boxed{
\begin{aligned}
&\sum_{j:u_j=0}z_j+\sum_{j:u_j=1}(1-z_j)\ge1,
&&u\in U,\\
&z_j\in\{0,1\},
&&j=1,\ldots,8.
\end{aligned}
}
\]
由于$u$是已知的常量向量，每条不等式也可以整理为标准线性形式
\[
\sum_{j:u_j=0}z_j-\sum_{j:u_j=1}z_j
\ge1-\sum_{j=1}^{8}u_j,
\qquad u\in U.
\]

例如，输入为$0$、输出为$1$对应的点为$u=0000\,0001$。对于本实验的置换S盒，该点既不属于差分合法集合，也不属于线性合法集合，其排除约束为
\[
z_1+z_2+z_3+z_4+z_5+z_6+z_7+(1-z_8)\ge1.
\]
将$0000\,0001$代入，左端为$0$；改变任意一位，左端便至少为$1$。这说明约束由非法点的坐标直接生成，不需要求解器先寻找其系数。

\paragraph{正确性分析}

对于固定的非法点$u$，若$z=u$，则$d_u(z)=0$，违反对应约束；若$z\ne u$，则至少一个坐标不同，从而$d_u(z)\ge1$。因此，一条约束在二进制空间中恰好排除其对应的非法点。

若$z\in V$，则$z$不同于每个$u\in U$，故满足全部不等式。若$z\notin V$，则$z\in U$，模型中存在针对$u=z$的约束，该约束必然被违反。由此可知，模型的二进制可行点集恰好为$V$。

\paragraph{用于差分与线性传播}

差分建模时，将DDT零位置对应的八位向量作为$U$；线性建模时，将LAT零位置对应的八位向量作为$U$。两种情况下均只使用原始变量$z$，不需要额外的选择变量。

\subsection{逻辑化简法}
\label{subsec:lab4-logic-minimization}

\paragraph{基本原理}

逐点排除法对每个非法点分别建立约束，而非法点之间可能共享部分坐标。逻辑化简法将这些点作为布尔函数的取值为$1$的输入，通过Espresso等工具合并为cube，再对每个cube建立一条排除约束。一条这样的约束可以同时排除该cube覆盖的全部非法点。

\paragraph{非法点的布尔表示}

设$U=\{0,1\}^{8}\setminus V$，定义非法点指示函数
\[
f_{\mathrm{bad}}(z)=
\begin{cases}
1,&z\in U,\\
0,&z\in V.
\end{cases}
\]
将该函数的真值表或积之和表达式交给逻辑化简工具，得到cube集合$\mathcal C$。每个cube写成
\[
c=(c_1,\ldots,c_8)\in\{0,1,-\}^{8},
\]
其中$0$、$1$表示该坐标固定，$-$表示该坐标不限。记
\[
I_0(c)=\{j\mid c_j=0\},
\qquad
I_1(c)=\{j\mid c_j=1\},
\]
则该cube覆盖的二进制点集为
\[
B(c)=\{z\in\{0,1\}^{8}\mid
z_j=0\ (j\in I_0(c)),\ z_j=1\ (j\in I_1(c))\}.
\]
精确化简要求覆盖关系满足
\begin{equation}
\bigcup_{c\in\mathcal C}B(c)=U.
\label{eq:lab4-cube-cover}
\end{equation}
该等式既要求所有非法点都被覆盖，也要求没有合法点落入任何cube。构造输入真值表时，合法点应明确标为$0$，不能将其当作无关项；cube中的$-$仅表示化简后该坐标不受限制。

讲义使用\texttt{espresso-logic}读取PLA格式的非法点列表：本实验应声明八个输入、一个输出，输入标签按$a_3,a_2,a_1,a_0,b_3,b_2,b_1,b_0$排列，每个非法点的输出标为$1$，未列出的点按对应PLA输入模式解释为输出$0$，不能当作任意取值。读取化简结果时，逐行提取输出为$1$的cube；其中的\texttt{-}表示该位不限，\texttt{.p}字段表示输出的cube数量。无论采用讲义中的\texttt{-Dexact}方式，还是通过其他接口调用Espresso，后续线性化都依赖式\eqref{eq:lab4-cube-cover}所要求的精确覆盖关系。

\paragraph{cube的线性约束}

为了排除$B(c)$，要求$z$至少有一个固定坐标不同于$c$，得到
\begin{equation}
\boxed{
\begin{aligned}
&\sum_{j\in I_0(c)}z_j
+\sum_{j\in I_1(c)}(1-z_j)\ge1,
&&c\in\mathcal C,\\
&z_j\in\{0,1\},
&&j=1,\ldots,8.
\end{aligned}
}
\label{eq:lab4-cube-model}
\end{equation}
整理常数项后，每条约束可写为
\[
\sum_{j\in I_0(c)}z_j-\sum_{j\in I_1(c)}z_j
\ge1-|I_1(c)|.
\]
cube中取值为$-$的坐标不出现在不等式中；若一个cube的全部坐标都固定，上式便退化为逐点排除约束。

\paragraph{正确性分析}

对于给定的$c$，若$z\in B(c)$，则所有固定坐标均与$c$一致，式\eqref{eq:lab4-cube-model}中对应排除不等式的左端为$0$，从而该点被排除。若$z\notin B(c)$，则至少一个固定坐标不同，左端至少为$1$，满足对应不等式。因此，每条不等式恰好排除一个cube覆盖的点集。

结合式\eqref{eq:lab4-cube-cover}可知，全部约束恰好排除$U$，留下的二进制可行点集就是$V$。用于差分传播时，对DDT零位置化简；用于线性传播时，对LAT零位置化简。化简后的约束只涉及原始变量，但启发式化简并不保证得到全局最少的cube，仍需核对覆盖关系是否精确。

\subsection{凸包法}
\label{subsec:lab4-convex-hull}

\paragraph{基本原理}

凸包法把合法的二进制点集$V$看作实数空间中的点集，先求出其凸包，再用凸包的线性不等式和等式描述传播关系。由于最终仍要求原始变量取二进制值，完整H表示在这些变量上的可行点集恰好为$V$。

\paragraph{凸包与H表示}

设$V=\{v^{(1)},\ldots,v^{(N)}\}\subseteq\{0,1\}^{8}$，其凸包定义为
\[
P=\operatorname{conv}(V)
=\left\{\sum_{i=1}^{N}\theta_i v^{(i)}\ \middle|\
\theta_i\ge0,\ \sum_{i=1}^{N}\theta_i=1\right\}.
\]
其中$\theta_i$为非负实数，此处仅用于给出凸包定义。将合法点作为顶点输入SageMath等工具，可得到$P$的H表示：
\[
P=\left\{z\in\mathbb R^{8}\ \middle|\
\begin{aligned}
&(a^{(k)})^{\mathsf T}z+\beta_k\ge0,
&&k\in K,\\
&(e^{(\ell)})^{\mathsf T}z+\gamma_\ell=0,
&&\ell\in E
\end{aligned}
\right\}.
\]
这里$K$、$E$分别为不等式和等式的索引集合，系数均为已知常数。当凸包不是满维时，H表示可能包含等式，读取结果时需要保留这些等式。为避免系数舍入改变可行点集，可使用精确有理数计算，并通过乘以正的公分母将约束整理成整数系数形式。

讲义通过\texttt{Polyhedron(vertices=V)}从点集构造多面体，并用\texttt{inequalities\_list()}读取系数。本实验每个点应为八位向量，因此输出的一行系数写成$q=(c_0,c_1,\ldots,c_8)$，对应
\[
h_q(z)=c_0+\sum_{j=1}^{8}c_jz_j\ge0.
\]
其中第一项$c_0$是常数项，不是第一个坐标的系数。等式系数需要另外读取，并以$h_q(z)=0$加入模型；不能仅提取不等式而丢弃等式条件。

由此得到完整凸包模型
\begin{equation}
\boxed{
\begin{aligned}
&(a^{(k)})^{\mathsf T}z+\beta_k\ge0,
&&k\in K,\\
&(e^{(\ell)})^{\mathsf T}z+\gamma_\ell=0,
&&\ell\in E,\\
&z_j\in\{0,1\},
&&j=1,\ldots,8.
\end{aligned}
}
\label{eq:lab4-full-hull-model}
\end{equation}

\paragraph{完整模型的正确性}

首先，任意$v\in V$都属于其凸包，故满足全部H表示约束。反过来，若二进制向量$z\in P$，则存在非负系数$\theta_i$，使$z=\sum_i\theta_i v^{(i)}$且$\sum_i\theta_i=1$。

当$z_j=0$时，所有$\theta_i>0$的点在第$j$个坐标都必须为$0$；当$z_j=1$时，这些点在第$j$个坐标都必须为$1$，否则加权平均会小于$1$。因此，所有正系数对应的点都与$z$逐坐标相同，必有$z\in V$。于是
\[
P\cap\{0,1\}^{8}=V,
\]
说明式\eqref{eq:lab4-full-hull-model}恰好描述合法点集。

\paragraph{基于非法点覆盖的贪心筛选}

完整H表示可能包含较多约束，而本实验只要求二进制可行点准确。因此，可以按照每条不等式排除非法点的能力进行筛选。设$U=\{0,1\}^{8}\setminus V$，对每个$k\in K$定义
\[
R_k=\{u\in U\mid (a^{(k)})^{\mathsf T}u+\beta_k<0\}.
\]
保留全部等式，并记等式已经排除的非法点集为
\[
R_E=\{u\in U\mid
\exists\ell\in E,\ (e^{(\ell)})^{\mathsf T}u+\gamma_\ell\ne0\}.
\]
初始化尚未排除的点集$U_0=U\setminus R_E$，选中的不等式索引集合$K_0=\varnothing$。只要$U_t\ne\varnothing$，便选择
\[
k_t\in\underset{k\in K\setminus K_t}{\operatorname{arg\,max}}
\ |R_k\cap U_t|,
\]
并更新
\[
K_{t+1}=K_t\cup\{k_t\},
\qquad
U_{t+1}=U_t\setminus R_{k_t}.
\]
若多个候选约束覆盖相同数量的点，可约定按固定的系数顺序选择，以便复现筛选过程。由于完整H表示能够排除每个非法二进制点，只要仍有点未被排除，就应存在能够覆盖其中至少一点的候选约束。

这对应讲义中的三个列表：\texttt{remaining}保存尚未排除的非法点，\texttt{pool}保存未选中的候选不等式，\texttt{chosen}保存已经保留的不等式。选中一条约束后，只从\texttt{remaining}中删除违反它的点，仍满足它的点继续留在列表中；同时将该约束从候选列表移至保留列表。若剩余点非空而最佳候选不能再排除任何点，应检查H表示、位序或系数读取，不能继续把该过程当作有效约简。

当剩余点集为空时，记最终选中的索引集合为$K_*$，保留的模型为
\[
\boxed{
\begin{aligned}
&(a^{(k)})^{\mathsf T}z+\beta_k\ge0,
&&k\in K_*,\\
&(e^{(\ell)})^{\mathsf T}z+\gamma_\ell=0,
&&\ell\in E,\\
&z_j\in\{0,1\},
&&j=1,\ldots,8.
\end{aligned}
}
\]
所有保留的约束均来自原H表示，因此不会排除合法点；筛选结束又保证每个非法点至少违反一个保留约束，故约简模型的二进制可行点集仍恰好为$V$。

这里的约简是相对于二进制点集而言的：约简后的连续可行区域可能大于原凸包，且贪心选择不保证约束数量全局最少。差分传播时取$V=V_{\mathrm D}$，线性传播时取$V=V_{\mathrm L}$，分别计算H表示并独立进行筛选。


\subsection{全解枚举与等价性验证}

按照讲义的验证方法，模型建立后将目标函数设为常数$0$，仅枚举满足约束的全部二进制解，不按概率或相关度大小筛选。Gurobi的\texttt{PoolSearchMode=2}用于系统搜索多个解，\texttt{PoolSolutions}指定希望收集的解数。原始变量只有八位，取值空间为$\{0,1\}^{8}$；对于上述四种模型，one-hot的选择变量由合法点唯一确定，其余模型不含辅助变量，因此模型的不同完整解也不超过原始八位点的取值总数。将解池容量设为严格大于这一上界，并在求解以\texttt{OPTIMAL}状态结束且解池未满时，才确认全部解已收集。超时或中断时只记录已有解，不据此声称枚举完成。

读取解池时，逐个设置\texttt{SolutionNumber}，再获取该解中各原始变量的值。讲义使用\texttt{Xn}属性，新版本Gurobi也提供\texttt{PoolNX}属性。对one-hot模型，还应核对恰有一个选择变量取$1$，且该选择变量对应的点与原始八位变量相同。

记枚举结果在原始变量上的投影集合为$S_{\mathrm{enum}}$，则分别计算
\[
V\setminus S_{\mathrm{enum}},
\qquad
S_{\mathrm{enum}}\setminus V.
\]
前者表示漏掉的合法点，后者表示多出的非法点。只有枚举已经完成且两个差集均为空，才说明模型的二进制解恰好为$V$。差分与线性两类点集需要分别进行验证；解数相同只是必要条件，还需逐点比较。实际约束系数、条数、枚举结果及求解状态在实验内容部分记录。
